\documentclass[12pt]{article}

% هوامش
\usepackage{geometry}
\geometry{margin=0.8in}

% ألوان/رؤوس/تذييل/تباعد
\usepackage{xcolor}
\usepackage{fancyhdr}
\usepackage{setspace}
\usepackage[inline]{enumitem}

% دعم يونيكود ولغات (XeLaTeX)
\usepackage{fontspec}
\usepackage{polyglossia}
\usepackage{bidi} % لدعم LTR/RTL

\usepackage{etoolbox} % لاختبار الفراغ

% Math
\usepackage{amsmath, amssymb, mathtools}

% (احتياط)
\usepackage{tabularx, ragged2e, array}
\newcolumntype{Y}{>{\RaggedRight\arraybackslash}X}

% عنوان السؤال كسطر مستقل يسار الصفحة (LTR)
% عنوان السؤال كسطر مستقل يسار الصفحة (LTR)
\newcommand{\qheader}[1]{%
  \par
  \begin{LTR}\noindent\textenglish{\textbf{Question~#1}}\end{LTR}%
  \par\vspace{0.2em}%
}

% خيارات سطرية مع التفاف تلقائي
% نُبقي المحتوى RTL (طبيعي للعربية)، ونجعل الوسم A., B. ... فقط LTR
\newcommand{\renderchoices}[5]{%
  \begin{enumerate*}[
    label=\protect\LR{\textenglish{\textbf{\Alph*.}}},
    itemjoin=\hspace{2em},
    itemsep=0.25em,
    leftmargin=0pt
  ]%
    \ifstrempty{#1}{}{\item #1}%
    \ifstrempty{#2}{}{\item #2}%
    \ifstrempty{#3}{}{\item #3}%
    \ifstrempty{#4}{}{\item #4}%
    \ifstrempty{#5}{}{\item #5}%
  \end{enumerate*}\par
}


\setmainlanguage[numerals=western]{arabic}
\setotherlanguage{english}

% الخطوط
\newfontfamily\arabicfont[Script=Arabic,Scale=1.05]{Amiri}
\newfontfamily\englishfont{Times New Roman}

% ترويسة/تذييل
\pagestyle{fancy}
\fancyhf{}
\renewcommand{\headrulewidth}{0pt}
\cfoot{\textenglish{Page } \thepage}

% فقرات
\setlength{\parskip}{0.4em}
\setlength{\parindent}{0pt}

\begin{document}
% ======= Header (fixed layout) =======
\noindent
\begin{tabular*}{\textwidth}{@{\extracolsep{\fill}} l c r}
\textenglish{Date: {\begin{LTR}{\textenglish{September 4, 2025}}\end{LTR}}} &
{\Large \textbf{{\textenglish{Algorithms -1-}}}} &
\textenglish{{\begin{LTR}{\textenglish{University Of Aleppo}}\end{LTR}}} \\
\textenglish{Student Name:} &
\textenglish{Model: {\begin{LTR}{\textenglish{C}}\end{LTR}}} &
\textenglish{\textbf{Duration:} {\begin{LTR}{\textenglish{90 minuts}}\end{LTR}}} \\
\textenglish{Student Number:} & &
\end{tabular*}

\vspace{0.9cm}

\section*{\begin{LTR}\textenglish{Select the correct answer for each question:}\end{LTR}}
\setstretch{1.25}

\noindent \qheader{1} \begin{LTR}{\textenglish{Find the inverse of the following matrix: }}\(\begin{bmatrix} 1 & 3 & 4 \\ 5 & 2 & 1 \\ 6 & 8 & 1 \end{bmatrix}\)\end{LTR}

\renderchoices{\begin{LTR}{\textenglish{There is no inverse for this matrix}}\end{LTR}}{\begin{LTR}\(\begin{pmatrix} 1 & 2 & 3 \\ 4 & 5 & 6 \\ 7 & 8 & 9  \end{pmatrix}\)\end{LTR}}{\begin{LTR}{\textenglish{0}}\end{LTR}}{\begin{LTR}{\textenglish{0}}\end{LTR}}{}
{\begin{LTR}\hfill\textenglish{[4 Points]}\end{LTR}}\par\vspace{0.5em}

\noindent \qheader{2} {\textarabic{أنت في مطعم وطلبت طبقاً. بعد أن تأكل الطبق، يأتي النادل ويضع طبقًا آخر فوقه. ما هي بنية البيانات التي تشبه هذا الموقف؟}}

\renderchoices{{\textarabic{مصفوفة (Array)}}}{{\textarabic{ مكدس (Stack)}}}{\begin{LTR}{\textenglish{0}}\end{LTR}}{\begin{LTR}{\textenglish{0}}\end{LTR}}{\begin{LTR}{\textenglish{0}}\end{LTR}}
{\begin{LTR}\hfill\textenglish{[1 Point]}\end{LTR}}\par\vspace{0.5em}

\noindent \qheader{3} \begin{LTR}\(x=3+6 \){\textenglish{ what is the value of x?}}\end{LTR}

\renderchoices{\begin{LTR}{\textenglish{7}}\end{LTR}}{\begin{LTR}{\textenglish{9}}\end{LTR}}{\begin{LTR}{\textenglish{8}}\end{LTR}}{\begin{LTR}{\textenglish{10}}\end{LTR}}{}
{\begin{LTR}\hfill\textenglish{[1 Point]}\end{LTR}}\par\vspace{0.5em}

\noindent \qheader{4} {\textarabic{ما هي الخوارزمية التي تستخدمها لتنظيم جوربك الملونة حسب اللون، ولكن فقط إذا كان لديك 3 جوارب؟}}

\renderchoices{{\textarabic{خوارزمية الفرز السريع (Quick Sort)}}}{{\textarabic{ خوارزمية الفرز بالإدراج (Insertion Sort)}}}{\begin{LTR}{\textenglish{0}}\end{LTR}}{\begin{LTR}{\textenglish{0}}\end{LTR}}{\begin{LTR}{\textenglish{0}}\end{LTR}}
{\begin{LTR}\hfill\textenglish{[5 Points]}\end{LTR}}\par\vspace{0.5em}

\noindent \qheader{5} {\textarabic{احسب مشتق التابع }}\( f(x) = x^2\){\textarabic{ وأوجد نهايته عند اللانهاية }}\(\lim_{x \to \infty}f(x)\)

\renderchoices{\begin{LTR}{\textenglish{14}}\end{LTR}}{\begin{LTR}{\textenglish{15}}\end{LTR}}{\begin{LTR}{\textenglish{16}}\end{LTR}}{\begin{LTR}{\textenglish{17}}\end{LTR}}{\begin{LTR}{\textenglish{17}}\end{LTR}}
{\begin{LTR}\hfill\textenglish{[1 Point]}\end{LTR}}\par\vspace{0.5em}

\noindent \qheader{6} {\textarabic{إذا كانت لدينا خوارزمية تستغرق نفس الوقت لحل مشكلة بغض النظر عن حجمها، فما هو اسمها؟}}

\renderchoices{\begin{LTR}{\textenglish{O(n)}}\end{LTR}}{\begin{LTR}{\textenglish{O(}}\(n^{2}\){\textenglish{)}}\end{LTR}}{\begin{LTR}{\textenglish{0}}\end{LTR}}{\begin{LTR}{\textenglish{0}}\end{LTR}}{}
{\begin{LTR}\hfill\textenglish{[1 Point]}\end{LTR}}\par\vspace{0.5em}

\noindent \qheader{7} {\textarabic{إذا كان لديك قائمة طويلة من أسماء الفواكه، لكنك تبحث عن "تفاحة". أي خوارزمية بحث ستستخدمها إذا كنت كسولاً وتتوقع أن "تفاحة" ستكون في مكان ما في منتصف القائمة؟}}

\renderchoices{{\textarabic{ البحث في الاتساع أولاً (Breadth-First Search)}}}{{\textarabic{ البحث الثنائي (Binary Search)}}}{\begin{LTR}{\textenglish{0}}\end{LTR}}{\begin{LTR}{\textenglish{0}}\end{LTR}}{\begin{LTR}{\textenglish{0}}\end{LTR}}
{\begin{LTR}\hfill\textenglish{[5 Points]}\end{LTR}}\par\vspace{0.5em}

\noindent \qheader{8} \begin{LTR}{\textenglish{Solve for }}\(x \){\textenglish{ in the equation }}\(2x+6 = 0 \)\end{LTR}

\renderchoices{\begin{LTR}{\textenglish{-1}}\end{LTR}}{\begin{LTR}{\textenglish{-2}}\end{LTR}}{\begin{LTR}{\textenglish{-3}}\end{LTR}}{\begin{LTR}{\textenglish{3}}\end{LTR}}{}
{\begin{LTR}\hfill\textenglish{[1 Point]}\end{LTR}}\par\vspace{0.5em}



\end{document}
